\section*{v1.7.0 Closure--Projection Synthesis Core}
\addcontentsline{toc}{section}{v1.7.0 Closure--Projection Synthesis Core}

\begin{remark}[Normalization note for v1.7.0]
\label{rem:normalization-note-v170}
In the present release, the mathematical content of this synthesis layer is subsumed by the OP transport theorem and the OP3 residual/shell interface. The block is retained as historical compression logic and interpretive guidance. This synthesis block is therefore kept as a historical transport precursor: its closure--projection synthesis is now carried proof-theoretically by Phase~II-D, Phase~II-E, and Phase~II-F, and here it should be read as an earlier integrative summary rather than as an independent axiomatic center.
\end{remark}

This release adds a further synthesis layer between the closure blocks already
introduced and the inherited structural core of the monolith. Its purpose is to
state, in one explicit place, how closure statements are transported across the
probe sector, how projection claims inherit their admissibility status from the
closure architecture, and how appendix tables are to be used as reversible proof
interfaces rather than as detached summaries.

\subsection*{1. Closure-to-projection transport discipline}
The central operational rule of v1.7.0 is that no projection statement may be
read as primary if its admissibility has not first been transported through the
closure chain. In other words, the order of explanation remains
\[
\begin{aligned}
&\text{local structure}
\Longrightarrow \text{HC}
\Longrightarrow \text{closure}\\[0.15em]
&\Longrightarrow \text{probe-compatible reduction}
\Longrightarrow \text{projection readout}.
\end{aligned}
\]
A projection is therefore not an independent object but a compressed exposure of
closure-stable content under a chosen observational or effective frame. This
rule keeps the paper aligned with the structure-first programme and prevents
later phenomenological language from silently replacing the closure mechanism
that generated it.

\paragraph{Synthesis Proposition A (closure status survives transport).}
If a statement has been established in the closure regime and is moved to a
probe-compatible or projective representation by a transport that preserves the
common body, then its structural status is not lost. What changes is the mode of
appearance, not the underlying admissibility class.

\paragraph{Proof sketch.}
The common body is fixed by the ground condition and stabilized by HC. Any
transport that preserves this body may change coordinates, phase labels, or
sector language, but it does not create a new ontology. Hence closure status is
inherited through the transport.

\subsection*{2. Probe reduction as a controlled interface}
The probe sector is not only a search device; in v1.7.0 it is also an
interface theorem. Two probes moving against one another, each carrying triolic
organization and coupled to an orthogonal third direction, define a reduced test
geometry in which admissible transitions can be filtered before full projection.
This means that the probe sector is a disciplined compression stage: it narrows
candidate structure while still remaining within closure control.

\paragraph{Synthesis Proposition B (probe reduction preserves closure discipline).}
Whenever the probe sector is constructed from triolic substructures under the
standing orthogonality and HC rules, the reduced search geometry cannot produce
admissible outputs that violate closure. It may discard candidates, rank them,
or expose only a subset of modes, but it cannot certify a structurally illegal
sector.

\paragraph{Proof sketch.}
The probe geometry is not external to the closure architecture. It is built from
substructures already governed by HC and orthogonality. Therefore any reduction
performed inside the probe channel is a reduction of representation, not a
release from confinement.

\subsection*{3. Projection tables as reversible proof interfaces}
The appendix tables introduced in the previous versions are to be read neither
as cosmetic summaries nor as final stand-alone derivations. Their exact role is
reversible: each table entry points upward to the closure and lemma chain from
which it is licensed, and downward to the sector language in which it may be
used. In this sense, a table row functions as a compact proof interface.

\paragraph{Synthesis Principle (table row admissibility).}
A table row is admissible only if one can move in both directions:
\begin{enumerate}[leftmargin=2em,label=(\alph*)]
  \item upward, from the row to the closure/lemma chain that stabilizes it; and
  \item downward, from the row to a projective or phenomenological reading that
  does not exceed the transported status of the original statement.
\end{enumerate}
If either direction is unavailable, the row is only heuristic and must be marked
as such.

\subsection*{4. Micro-to-macro synthesis route}
The micro-to-macro route is now stated in a more disciplined form. The local
triolic unit is not yet a macroscopic law, but it already fixes the admissible
coupling pattern. HC then stabilizes multiple appearances of the same body,
closure confines the admissible dynamics, the probe sector performs controlled
reduction, and only then do effective macroscopic equations appear as readouts.
The macro level is therefore a closure-saturated exposure of the micro level,
not a rival explanatory layer.

\paragraph{Synthesis Proposition C (macro equations are closure-saturated readouts).}
Under the standing assumptions of the monolith, any effective macroscopic
relation used later in the paper must be readable as a closure-saturated image
of the local structural chain. Hence macro equations are legitimate exactly to
the degree that the closure-to-projection transport remains traceable.

\subsection*{5. Appendix back-reference discipline, strengthened}
The appendix is now explicitly promoted to a back-reference layer. This means
that every compressed proof move used in the running text should be recoverable
through one of three routes: a lemma reference, a tool reference, or a table
reference with reversible upward justification. The appendix is therefore not an
archive placed after the fact; it is the controlled external memory of the proof
architecture.

\paragraph{Operational rule for later versions.}
Whenever a derivation step is abbreviated in the main line, the abbreviated step
should be tagged by a reusable interface name that can be expanded again from
the appendix without changing the logical order of the argument. This preserves
speed without sacrificing structural auditability.


